% SGA 3, Exposé IX, tradução brasileira-portuguesa.
% Escopo: §7 integral.
% Autoridade: Polo--Gille Exp9-8nov09.pdf, pp. locais 26--30.

\SGARefTarget{sga3:ix:section:7}
\section*{7. Algebricidade dos homomorfismos formais em um grupo afim}
\addcontentsline{toc}{section}{7. Algebricidade dos homomorfismos
formais em um grupo afim}

\medskip
\noindent\SGARefTarget{sga3:ix:theorem:7:1}\textbf{Teorema 7.1.}
\textit{Seja \(A\) um anel noetheriano munido de um ideal \(I\) tal
que \(A\) seja separado e completo para a topologia \(I\)-ádica. Ponhamos
\[
  S=\operatorname{Spec}(A),
  \qquad
  S_n=\operatorname{Spec}(A/I^{n+1}),
\]
e sejam \(H,G\) pré-esquemas em \(S\)-grupos, com \(H\) de tipo
multiplicativo isotrivial e \(G\) afim. Então, a aplicação canônica
\SGARefTarget{sga3:ix:theorem:7:1:equation:x}
\begin{equation}
  \tag{x}
  \theta:
  \operatorname{Hom}_{S\text{-}\mathrm{gr}}(H,G)
  \longrightarrow
  \varprojlim_n
  \operatorname{Hom}_{S_n\text{-}\mathrm{gr}}(H_n,G_n)
\end{equation}
é bijetiva, onde \(H_n,G_n\) são os \(S_n\)-grupos deduzidos de
\(H,G\) pela mudança de base \(S_n\to S\).}

\emph{Demonstração.}
Suponhamos primeiro que \(H\) é diagonalizável, portanto da forma
\[
  H=\operatorname{Spec}A(M),
\]
onde \(B=A(M)\) é a álgebra do grupo comutativo \(M\) com
coeficientes em \(A\). Podemos escrever também
\[
  G=\operatorname{Spec}(C),
\]
onde \(C\) é uma \(A\)-álgebra munida de uma aplicação diagonal que
satisfaz os axiomas usuais. Os homomorfismos de \(S\)-grupos
\(H\to G\) correspondem bijetivamente aos homomorfismos de
\(A\)-álgebras
\[
  \varphi:C\longrightarrow B
\]
compatíveis com as aplicações diagonais, isto é, tais que, para
todo \(f\in C\),
\SGARefTarget{sga3:ix:theorem:7:1:display:01}
\[
  \Delta_H(\varphi(f))
  =(\varphi\otimes\varphi)(\Delta_G(f)).
\]
Uma descrição análoga vale para os homomorfismos de
\(S_n\)-grupos \(H_n\to G_n\), dados por homomorfismos de
\(A_n\)-álgebras
\[
  \varphi_n:C_n\longrightarrow B_n,
\]
onde
\[
  A_n=A/I^{n+1},
  \qquad
  B_n=B\otimes_AA_n,
  \qquad
  C_n=C\otimes_AA_n.
\]
\emph{[A fonte imprime \(C=C\otimes_AA_n\) na última igualdade;
entende-se \(C_n\).]}
Ponhamos
\SGARefTarget{sga3:ix:theorem:7:1:display:02}
\[
  \widehat B=\varprojlim_n B_n,
  \qquad
  \widehat C=\varprojlim_n C_n.
\]
Então,
\[
  \widehat B=\varprojlim_n A_n(M)
\]
identifica-se com um submódulo do produto \(A^M\), a saber, o
submódulo formado pelas famílias \((a_m)_{m\in M}\) de elementos de
\(A\) que tendem a zero para a topologia \(I\)-ádica segundo o filtro
dos complementares das partes finitas de \(M\). Em particular, o
homomorfismo canônico
\[
  B\longrightarrow\widehat B
\]
é injetivo.\footnote{Nota do editor: desenvolveu-se o argumento que
segue.}
O sistema projetivo
\[
  \theta(\varphi)=(\varphi_n)_{n\geq0}
\]
define um homomorfismo
\[
  \widehat\varphi:\widehat C\longrightarrow\widehat B
\]
tal que \SGARefLink{sga3:ix:theorem:7:1:diagram:01}{o diagrama seguinte}
é comutativo:
\SGARefTarget{sga3:ix:theorem:7:1:diagram:01}
\[
\begin{tikzcd}[column sep=large,row sep=large]
  C
    \arrow[r,"\varphi"]
    \arrow[d]
  & B
    \arrow[d]
  \\
  \widehat C
    \arrow[r,"\widehat\varphi"]
  & \widehat B.
\end{tikzcd}
\]

Provemos que \(\theta\) é sobrejetiva. Seja
\((\varphi_n)_{n\geq0}\) um sistema projetivo; devemos mostrar que ele
é obtido por redução a partir de algum \(\varphi\) do primeiro membro de
\SGARefLink{sga3:ix:theorem:7:1:equation:x}{\((x)\)}. A priori, o
sistema \((\varphi_n)\) define um homomorfismo sobre as álgebras
completadas,
\[
  \widehat\varphi:\widehat C\longrightarrow\widehat B.
\]
Basta provar que a composta
\[
  \Phi:
  C\longrightarrow\widehat C
   \xrightarrow{\widehat\varphi}\widehat B
\]
envia \(C\) em \(B\). Com efeito, nesse caso ela fornece um
homomorfismo de \(A\)-álgebras
\SGARefTarget{sga3:ix:theorem:7:1:display:phi:01}\[
  \varphi:C\longrightarrow B
\]
cujas reduções são os \(\varphi_n\). Vê-se imediatamente que
\SGARefLink{sga3:ix:theorem:7:1:display:phi:01}{esse homomorfismo} é
compatível com as aplicações diagonais, já que os
\(\varphi_n\) o são e uma vez que
\[
  B\otimes_AB\longrightarrow\widehat{B\otimes_AB}
\]
é injetivo, como resulta do mesmo argumento após substituir \(M\)
por \(M\times M\).

Ao passar aos completamentos, o homomorfismo diagonal \(\Delta_H\)
define
\SGARefTarget{sga3:ix:theorem:7:1:display:03}
\[
  \widehat\Delta_H:
  \widehat B\longrightarrow\widehat{B\otimes_AB}.
\]
Tomando o limite projetivo dos diagramas análogos definidos pelos
\(\varphi_n\), obtém-se o diagrama comutativo
\SGARefTarget{sga3:ix:theorem:7:1:diagram:02}
\[
\begin{tikzcd}[column sep=huge,row sep=large]
  C
    \arrow[r,"\Phi"]
    \arrow[d,"\Delta_G"']
  & \widehat B
    \arrow[d,"\widehat\Delta_H"]
  \\
  C\otimes_A C
    \arrow[r,"\Psi"]
  & \widehat{B\otimes_A B},
\end{tikzcd}
\]
onde \(\Psi\) é a composta
\SGARefTarget{sga3:ix:theorem:7:1:display:04}
\[
  C\otimes_A C
  \xrightarrow{\ \Phi\otimes\Phi\ }
  \widehat B\otimes_A\widehat B
  \longrightarrow
  \widehat{B\otimes_A B};
\]
A última flecha é o homomorfismo canônico evidente. Assim,
para todo \(f\in C\), o elemento \(\Phi(f)\in\widehat B\) tem a
propriedade de que sua imagem por \(\widehat\Delta_H\) é um elemento
«decomponível» de \(\widehat{B\otimes_A B}\), a saber, um elemento da
imagem de \(\widehat B\otimes_A\widehat B\).

Seja \((e_m)_{m\in M}\) a base canônica de \(A(M)\), e seja
\((e_{m,m'})\) a de
\[
  A(M\times M)=A(M)\otimes_AA(M).
\]
Como
\[
  \Delta_H(e_m)=e_m\otimes e_m=e_{m,m}
  \qquad (m\in M),
\]
basta agora aplicar \SGARefLink{sga3:ix:lemma:7:2}{o lema seguinte}.

\medskip
\noindent\SGARefTarget{sga3:ix:lemma:7:2}\textbf{Lema 7.2.}
\begin{itshape}
Sejam \(A\) um anel noetheriano, \(M\) um conjunto e
\((a_{m,m'})\) uma família de elementos de \(A\), indexada por
\(M\times M\), que satisfaz:
\begin{enumerate}
  \item[(i)] \SGARefTarget{sga3:ix:lemma:7:2:item:i}
  \(a_{m,m'}=0\) se \(m\neq m'\); isto é, o suporte da família
  está contido na diagonal de \(M\times M\);

  \item[(ii)] \SGARefTarget{sga3:ix:lemma:7:2:item:ii}
  Tem-se
  \SGARefTarget{sga3:ix:lemma:7:2:display:01}
  \[
    a_{m,m'}=\sum_i b_m^ic_{m'}^i,
  \]
  onde os \(b^i,c^i\) são um número finito de elementos de \(A^M\);
  dito de outro modo, \((a_{m,m'})\) pertence à imagem do
  homomorfismo canônico
  \[
    A^M\otimes_AA^M\longrightarrow A^{M\times M}.
  \]
\end{enumerate}
Nessas condições, o suporte de \((a_{m,m'})\) é finito.
\end{itshape}

\emph{Demonstração.}
Por \SGARefLink{sga3:ix:lemma:7:2:item:i}{(i)}, a família
\((a_{m,m'})\) é determinada pelos elementos
\[
  a_m=a_{m,m}.
\]
Para \(x=(x_n)_{n\in M}\in A(M)\), ponhamos
\[
  (u\cdot x)_m=\sum_{m'}a_{m,m'}x_{m'}.
\]
Interpretamos, assim, \((a_{m,m'})\) como a matriz de um homomorfismo
\[
  u:A(M)\longrightarrow A^M.
\]
A condição \SGARefLink{sga3:ix:lemma:7:2:item:i}{(i)} dá simplesmente
\[
  (u\cdot x)_m=a_mx_m.
\]
Além disso, a condição
\SGARefLink{sga3:ix:lemma:7:2:item:ii}{(ii)} dá
\[
  u\cdot x\in\sum_i A\cdot b^i.
\]
Portanto, \(u(A(M))\) está contido em um \(A\)-módulo de tipo
finito. Se \((e_m)_{m\in M}\) designa a base canônica de
\[
  A(M)\subset A^M,
\]
todos os elementos \(a_me_m\) estão contidos em um \(A\)-módulo de
tipo finito. Como \(A\) é noetheriano, o módulo gerado por
esses elementos é, por sua vez, de tipo finito. Como os \(e_m\) são
linearmente independentes, todos os \(a_m\), salvo um número finito,
são nulos. Isso prova o lema e, portanto, o Teorema
\SGARefLink{sga3:ix:theorem:7:1}{7.1}, quando \(H\) é diagonalizável.
\qed

Tratemos agora o caso geral de
\SGARefLink{sga3:ix:theorem:7:1}{7.1}, no qual só se supõe que
\(H\) seja isotrivial. Existe, então, um morfismo finito étale
e sobrejetivo
\[
  S'\longrightarrow S
\]
tal que
\[
  H'=H\times_SS'
\]
é diagonalizável. Utilizaremos somente o fato de que \(S'\to S\)
é finito e de recobrimento para a topologia fielmente plana e
quase compacta, ou simplesmente para a topologia canônica sobre
\((\mathbf{Sch})\). Assim, «étale» poderia ser substituído por «plano».

Ponhamos
\[
  S''=S'\times_SS',
\]
e introduzamos também
\[
  S'_n
  \quad\text{e}\quad
  S''_n
  =S''\times_SS_n
  =S'_n\times_{S_n}S'_n,
\]
juntamente com \(H',G',H'',G'',H'_n,G'_n,H''_n,G''_n\), deduzidos de
\(H,G\) pelas mudanças de base evidentes. Os grupos \(H'\) e \(H''\)
são diagonalizáveis. Ademais, \(S'\), e, portanto, \(S''\), é afim. Se
\[
  S'=\operatorname{Spec}(A'),
  \qquad
  S''=\operatorname{Spec}(A''),
\]
então \(A'\) e \(A''\) são separados e completos para as topologias
definidas por \(IA'\) e \(IA''\), respectivamente, porque ambos são
finitos sobre \(A\).

Como \(S'\to S\) e \(S'_n\to S_n\) são morfismos de recobrimento,
temos um diagrama comutativo de aplicações de conjuntos cujas duas
linhas são exatas:
\SGARefTarget{sga3:ix:theorem:7:1:diagram:03}
\[
\begin{tikzcd}[
  column sep=small,
  row sep=large,
  cells={nodes={font=\small}}
]
  \operatorname{Hom}_{S\text{-}\mathrm{gr}}(H,G)
    \arrow[r]
    \arrow[d,"u"]
  & \operatorname{Hom}_{S'\text{-}\mathrm{gr}}(H',G')
    \arrow[r,shift left=.6ex]
    \arrow[r,shift right=.6ex]
    \arrow[d,"u'"]
  & \operatorname{Hom}_{S''\text{-}\mathrm{gr}}(H'',G'')
    \arrow[d,"u''"]
  \\
  \varprojlim_n
    \operatorname{Hom}_{S_n\text{-}\mathrm{gr}}(H_n,G_n)
    \arrow[r]
  & \varprojlim_n
    \operatorname{Hom}_{S'_n\text{-}\mathrm{gr}}(H'_n,G'_n)
    \arrow[r,shift left=.6ex]
    \arrow[r,shift right=.6ex]
  & \varprojlim_n
    \operatorname{Hom}_{S''_n\text{-}\mathrm{gr}}(H''_n,G''_n).
\end{tikzcd}
\]
A segunda linha é exata porque é o limite projetivo dos
diagramas exatos associados aos diferentes \(S'_n\to S_n\). Pelo
caso diagonalizável já provado, \(u'\) e \(u''\) são bijetivos. Por
conseguinte, \(u\) também o é, o que termina a demonstração de
\SGARefLink{sga3:ix:theorem:7:1}{7.1}. \qed

\medskip
\noindent\SGARefTarget{sga3:ix:corollary:7:3}\textbf{Corolário 7.3.}
\textit{Nas condições de
\SGARefLink{sga3:ix:theorem:7:1}{7.1}, suponhamos, ademais, que \(G\) seja
liso sobre \(S\), e seja
\[
  u_0:H_0\longrightarrow G_0
\]
um homomorfismo de \(S_0\)-grupos. Existe, então, um homomorfismo de
\(S\)-grupos
\[
  u:H\longrightarrow G
\]
que levanta \(u_0\). Quaisquer dois levantamentos \(u,u'\) são
conjugados por um elemento \(g\in G(S)\) cuja redução é o elemento
unidade de \(G_0(S_0)\).}

\emph{Demonstração.}
Isso resulta da combinação de
\SGARefLink{sga3:ix:theorem:3:6}{3.6} e
\SGARefLink{sga3:ix:theorem:7:1}{7.1}. Para construir \(u\), constroem-se
sucessivamente os \(u_n\), o que é possível por
\SGARefLink{sga3:ix:theorem:3:6}{3.6}; o Teorema
\SGARefLink{sga3:ix:theorem:7:1}{7.1} mostra, então, que o sistema
\((u_n)\) provém de um \(S\)-homomorfismo \(u\). Dados dois
levantamentos \(u,u'\), constroem-se sucessivamente elementos \(g_n\)
tais que
\[
  g_0=1,
  \qquad
  g_n\text{ é a redução de }g_{n+1},
  \qquad
  u'_n=\operatorname{int}(g_n)u_n.
\]
Mais uma vez, isso é possível por
\SGARefLink{sga3:ix:theorem:3:6}{3.6}. Como \(A\) é separado e
completo, os \(g_n\) provêm de um elemento \(g\in G(S)\). Por fim, a
injetividade de \SGARefLink{sga3:ix:theorem:7:1}{7.1} dá
\[
  u'=\operatorname{int}(g)u.
\]
\qed

\medskip
\noindent\SGARefTarget{sga3:ix:remark:7:4}\textbf{Observação 7.4.}
\textit{Compare-se \SGARefLink{sga3:ix:theorem:7:1}{7.1} com EGA III,
5.4.1, que implica que o enunciado de
\SGARefLink{sga3:ix:theorem:7:1}{7.1} continua válido se, em vez
de supor que \(H\) seja de tipo multiplicativo e que \(G\) seja afim,
supõe-se que \(H\) seja próprio sobre \(S\), e que \(G\) seja separado e
localmente de tipo finito sobre \(S\). Um enunciado como o de
\SGARefLink{sga3:ix:theorem:7:1}{7.1} sem hipótese de propriedade é
impossível em geral; ele é um dos aspectos mais importantes da
grande rigidez da estrutura de um grupo de tipo
multiplicativo.}

\textit{O enunciado análogo com
\[
  G=H=\mathbf G_a
\]
é falso em geral. Com efeito, tome-se \(A\) de característica \(p>0\),
e definam-se os \(u_n\), por redução módulo \(I^{n+1}\), a partir de
uma série formal aditiva
\SGARefTarget{sga3:ix:remark:7:4:display:01}
\[
  u(T)=\sum_{n\geq0}a_nT^{p^n},
\]
onde os \(a_n\in A\) tendem a zero para a topologia \(I\)-ádica,
mas não para a topologia discreta.}

\textit{O enunciado de
\SGARefLink{sga3:ix:theorem:7:1}{7.1} também se torna falso se se
suprime a hipótese de que \(G\) seja afim, mesmo quando
\(H=\mathbf G_m\). Obtém-se um exemplo, com \(A\) um anel de
valorização discreta completo, a partir de uma curva elíptica sobre o
corpo de frações \(K\) de \(A\) cuja redução, na teoria de
Néron--Kodaira, digamos, é o grupo \(\mathbf G_m\) sobre o corpo
residual \(k\).\footnote{Nota do editor: poder-se-ia dar explicitamente
um exemplo semelhante.} Tem-se, então, um esquema em grupos
comutativo liso \(G\) sobre \(S\), cuja fibra especial é
\(\mathbf G_{m,k}\). Por
\SGARefLink{sga3:ix:theorem:3:6}{3.6}, isso permite construir um sistema
projetivo de isomorfismos
\[
  u_n:H_n\xrightarrow{\ \sim\ }G_n,
  \qquad H=\mathbf G_{m,A},
\]
ao passo que a fibra genérica de \(G\) é uma variedade abeliana. Por
conseguinte, não existe nenhum homomorfismo de \(S\)-grupos \(H\to G\)
distinto do homomorfismo zero.}
